% Generated draft: requires mathematical and source review.
\documentclass[11pt]{article}
\usepackage[margin=25mm]{geometry}
\usepackage{fontspec}
\setmainfont{Latin Modern Roman}
\usepackage{amsmath,amssymb,mathtools,unicode-math}
\setmathfont{Latin Modern Math}
\usepackage{enumitem,needspace,xurl,hyperref}
\setlength{\parindent}{0pt}
\setlength{\parskip}{5pt}
\setlength{\emergencystretch}{3em}
\newcommand{\N}{\mathbb{N}}
\newcommand{\Z}{\mathbb{Z}}
\newcommand{\Q}{\mathbb{Q}}
\newcommand{\R}{\mathbb{R}}
\newcommand{\C}{\mathbb{C}}
\newcommand{\F}{\mathbb{F}}
\newcommand{\A}{\mathbb{A}}
\newcommand{\PP}{\mathbb{P}}
\newcommand{\E}{\mathbb{E}}
\newcommand{\eps}{\varepsilon}
\newcommand{\dd}{\,\mathrm{d}}
\begin{document}
\section*{UnsolvedMath ID: 30006175}
\textbf{Source problem code:} catalogue entry
\subsection{Definitions and mathematical statement}
A stochastic process is a family \(X=(X_t)_{t\in T}\) of random variables on a probability space, taking values in a measurable state space. The supplied record does not specify the index set \(T\), state space, probability spaces, or any path regularity assumptions; these are placeholders, not recovered hypotheses.

For a specified set \(\mathcal S\) of stochastic processes or their chosen representations, a metric is a function
\[
d:\mathcal S\times\mathcal S\longrightarrow[0,\infty)
\]
satisfying, for every \(X,Y,Z\in\mathcal S\),
\[
\begin{gathered}
d(X,Y)=0\iff X=Y,\qquad d(X,Y)=d(Y,X),\\
d(X,Z)\le d(X,Y)+d(Y,Z).
\end{gathered}
\]
Here equality depends on the chosen objects: equality of representations, equality in law, and a behavioral equivalence such as bisimulation are different possibilities. If distinct representations may have distance zero, the corresponding construction is a pseudometric on representations and can be a metric on a suitable quotient. The record does not select an equivalence relation.

\textbf{Source question.} How can one define similarity metrics between stochastic processes? The associated contribution is titled \emph{Optimal transport distances for Markov chains} and is attributed to Gergely Neu, jointly with Sergio Calo, Anders Jonsson, Ludovic Schwartz, and Javier Segovia-Aguas.

\textbf{Unresolved formulation.} The question asks for a meaningful notion of similarity, but supplies neither a particular class \(\mathcal S\), a criterion for similarity, nor additional requirements on \(d\). The metric axioms above explain the terminology; they do not reconstruct a specific existence theorem or a uniquely determined research target. The Markov-chain setting is indicated by the contribution title, but transition kernels, state-space assumptions, initial distributions, time horizons, transport costs, coupling restrictions, and any discount parameters are absent. No particular optimal-transport or bisimulation formula can therefore be stated from this record alone.

The supplied statement-verification note reports that the official report poses the broad question, presents an optimal-transport construction, identifies the relevant bisimulation metrics as optimal-transport distances, and gives a finite-dimensional algorithm for Markov chains. These are supplied source descriptions, not independently verified theorems here; their hypotheses and precise conclusions are not provided. The registry marks the statement as \texttt{unrecoverable} and its literature status as \texttt{partially\_solved}. Neither label resolves the missing mathematical specifications.
\subsection{Short English statement}
The source asks how to measure similarity between stochastic processes, with an associated contribution on optimal-transport distances for Markov chains. A precise target requires the process class, the meaning of equality or equivalence, and the desired similarity properties. These are missing, so this draft preserves the broad question and records the unresolved qualifications. The supplied partial-solution assessment requires independent review.
\subsection{Sources}
\begin{itemize}
\item Local UnsolvedMath registry, supplied complete record for permanent ID \(30006175\), title \emph{Similarity Metrics for Stochastic Processes}, source code \texttt{catalogue entry}; collection: Oberwolfach Reports Open Problems, proposed year 2025. No separate proposer is supplied.
\item \emph{Overparametrization, Regularization, Identifiability and Uncertainty in Machine Learning} (2025), Oberwolfach Reports, pp.\ 175--208, DOI \texttt{10.4171/owr/2025/4}. Supplied report link: \url{https://doi.org/10.4171/owr/2025/4}. The extracted statement includes page 197 and the contribution title \emph{Optimal transport distances for Markov chains}, attributed to Gergely Neu, jointly with Sergio Calo, Anders Jonsson, Ludovic Schwartz, and Javier Segovia-Aguas.
\item Supplied literature source: \emph{Oberwolfach Report 4/2025}, \url{https://content.ems.press/assets/public/full-texts/serials/owr/22/1/14299083/online/10.4171-owr-2025-4.pdf}.
\item Supplied literature source: Calo et al., \emph{Bisimulation metrics are optimal transport distances}, \url{https://arxiv.org/abs/2406.04056}.
\item Supplied literature source: Calo et al., \emph{Distances for Markov chains from sample streams}, \url{https://arxiv.org/abs/2505.18005}.
\item Supplied literature assessment dated \texttt{2026-08-22}: status \texttt{partially\_solved}. The supplied assessment states that optimal-transport and bisimulation constructions provide principled similarity distances for Markov chains, including efficient computation, while no universal metric theory for all stochastic processes is established. This is supplied dated curation derived from the OWR record, not an independently verified literature assessment. Time-sensitive status and cited claims require independent verification.
\end{itemize}
\end{document}
